\section{Conclusion}
\label{sec:conclusion}

We introduced the Bayesian X-Learner, a meta-learner that pairs
cross-fitted doubly robust pseudo-outcomes with a Welsch redescending
pseudo-likelihood to produce calibrated MCMC posteriors over $\tx$
under heavy-tailed outcomes. On the clean-data IHDP benchmark the
default configuration leads on mean $\PEHE$ across S-/T-/X-Learners,
Causal BART, and EconML's causal forest ($\PEHE = 0.56$, tightest
dispersion among competitive entries); on contaminated whale DGPs up to
$\sim$20--25\% density, a single-flag
\texttt{contamination\_severity="severe"} switch recovers RMSE
$\approx 0.13$ with calibrated 95\% intervals. We validated on the
Hillstrom email-marketing RCT ($N = 42{,}613$) and reported
covariate-stratified $\tau(x)$ coverage to substantiate calibration
beyond scalar summaries. We separated
tails-as-contamination (handled by Welsch + Huber nuisance) from
tails-as-signal (handled by a tail-aware CATE basis) and showed
empirically that conflating the two via a data-layer EVT rescaling
is actively harmful. Extended MCMC diagnostics (BFMI,
autocorrelation, initialisation sensitivity) confirm reliable
sampler behaviour for the non-convex Welsch posterior; a
strengthened $\eta$-calibration analysis with empirical verification
of the positive-definiteness condition provides firmer ground for
the generalised-Bayes credible intervals. Six empirical
boundaries define where the library is validated; outside them, we
point users to tools that are better suited. The contribution is
not a claim of state-of-the-art dominance but a principled closing
of the three-way gap among heterogeneity, calibration, and
heavy-tail robustness that practitioners actually face.

